% =============================================================================
\chapter{Worked Examples}
\label{ch:examples}
% =============================================================================

% -----------------------------------------------------------------------------
\section{Example 1: Curved Beam Optimisation}
\label{sec:ex:beam}
% -----------------------------------------------------------------------------

This example corresponds to \texttt{demo\_beam\_optimization.m}.  An S-shaped
beam is imported from an IGES file, clamped at one end and guided at the
other, then thickness-optimised.

\subsection{Problem Setup}

\begin{lstlisting}[caption={Material and loading definition.}]
E         = 69e9;       % Ti-6Al-4V, Young's modulus (Pa)
ultimate  = 900e6;      % Ultimate tensile strength (Pa)
rho       = 4820;       % Density (kg/m^3)
beam_width = 0.008;     % 8 mm out-of-plane width
initial_thickness = 0.003;  % 3 mm starting thickness
Fx = 0;  Fy = -333;    % 333 N vertical load at free end
\end{lstlisting}

\subsection{Geometry Import}

The curve is imported from a SolidWorks-exported \texttt{.igs} file and
resampled to 250 uniform-arc-length elements:

\begin{lstlisting}[caption={Importing the beam centreline.}]
import_opts = struct('num_elements', 250, 'scale', 0.001);
result = importIGSCurve('s_test_2.igs', import_opts);
\end{lstlisting}

\subsection{Boundary Conditions}

Node~1 is fully clamped (all three DOFs fixed).  The free end (last node) is
constrained in $x$ and $\theta$ but free in $y$ (a ``guided'' support):

\begin{lstlisting}[caption={Boundary conditions.}]
constraints = [1,         1, 0;    % ux = 0
               1,         2, 0;    % uy = 0
               1,         3, 0;    % theta = 0
               num_nodes, 1, 0;    % guided: ux = 0
               num_nodes, 3, 0];   % guided: theta = 0
\end{lstlisting}

\subsection{Analysis and Optimisation}

A nonlinear analysis is run first to establish baseline deflection and
stress.  The optimiser then iterates with a safety factor of~2:

\begin{lstlisting}[caption={Optimisation call.}]
opt = struct('alpha', 0.4, 'safety_factor', 2, ...
             'use_nonlinear', true, 't_min', 0.0005);
[geom, res] = optimizeThickness(geom, analysis, load_case, opt);
\end{lstlisting}

\subsection{Typical Results}

\begin{table}[H]
\centering
\caption{Beam optimisation results (representative).}
\begin{tabular}{lcc}
\toprule
\textbf{Metric} & \textbf{Initial} & \textbf{Optimised} \\
\midrule
Mass                    & 7--10~g    & 3--5~g \\
Max displacement        & 0.5--2~mm  & 1--4~mm \\
Min factor of safety    & $\gg 2$    & $\approx 2$ \\
Max stress              & 50--100~MPa & 400--450~MPa \\
Thickness range         & 3~mm (uniform) & 0.5--8~mm \\
\bottomrule
\end{tabular}
\end{table}

The optimiser redistributes material from low-stress regions (near the
neutral axis of curves) to high-stress regions (the outer fibres of bends),
producing a variable-thickness organic profile.

% -----------------------------------------------------------------------------
\section{Example 2: Truss Topology Optimisation}
\label{sec:ex:truss}
% -----------------------------------------------------------------------------

This example corresponds to \texttt{demo\_truss\_optimization.m}.  A Warren
truss is defined programmatically, subdivided into fine elements, and
optimised.

\subsection{Problem Setup}

A 6-panel Warren truss with 300~mm span and 100~mm depth:

\begin{lstlisting}[caption={Truss parameters.}]
span       = 0.300;       % 300 mm
height     = 0.100;       % 100 mm
n_panels   = 6;           % 7 bottom nodes, 6 top nodes
E          = 200e9;       % Steel
ultimate   = 400e6;
beam_width = 0.010;       % 10 mm width
initial_thickness = 0.003;
Fy_load    = -5000;       % 5 kN centre point load
\end{lstlisting}

\subsection{Truss Geometry}

Bottom-chord nodes are spaced uniformly at $y = 0$.  Top-chord nodes are at
half-panel offsets at $y = h$:

\begin{lstlisting}[caption={Warren truss node placement.}]
bot_x = (0:n_panels)' * dx_panel;
top_x = ((0:n_panels-1)' + 0.5) * dx_panel;
\end{lstlisting}

Connectivity includes bottom chord, top chord, and diagonal members forming
the characteristic alternating-V Warren pattern.

\subsection{Subdivision}

Each of the original members is split into 15 sub-elements:

\begin{lstlisting}[caption={Truss subdivision.}]
[sub_nodes, sub_conn, member_map] = subdivideTruss(...
    truss_nodes, conn, 15);
\end{lstlisting}

This creates approximately 320~nodes and 330~elements from the original
13~nodes and 22~members.

\subsection{Optimisation}

The optimiser uses a low minimum thickness to allow elements to nearly
vanish:

\begin{lstlisting}[caption={Truss optimisation settings.}]
opt = struct('alpha', 0.3, 'safety_factor', 2, ...
    't_min', 0.0002, 'max_iterations', 300);
[geom, res] = optimizeThickness(geom, analysis, load_case, opt);
\end{lstlisting}

\subsection{Interpreting Results}

After optimisation, material concentrates along the primary load paths
(top and bottom chords, and the most heavily loaded diagonals) while
non-structural diagonals thin to near-zero.  The result resembles a
continuous topology-optimised shape rather than a discrete truss.

The \texttt{min\_thickness\_frac} option in \texttt{exportProfile} filters
out vanished elements, leaving clean boundaries for CAD export:

\begin{lstlisting}[caption={Exporting the truss profile.}]
opts = struct('min_thickness_frac', 0.15, ...
    'export_filename', 'profiles/optimised_truss');
profile = exportProfile(geom, opts);
\end{lstlisting}

% -----------------------------------------------------------------------------
\section{Example 3: Minimal Custom Problem}
\label{sec:ex:custom}
% -----------------------------------------------------------------------------

A three-element L-bracket, built entirely from code with no file imports:

\begin{lstlisting}[caption={Complete minimal example.}]
% Geometry
nodes = [0 0; 0.1 0; 0.1 0.08];
conn  = [1 2; 2 3];

props = struct('E', 200e9, 'thickness', 0.005, ...
    'width', 0.012, 'ultimate', 350e6, 'rho', 7850);

constraints = [1 1 0; 1 2 0; 1 3 0];  % clamp node 1

geom     = StructureGeometry(nodes, conn, props, constraints);
analysis = FEAAnalysis(geom);

% Load: 2 kN downward at the tip
lc = LoadCase(geom);
lc.addNodalLoad(3, 0, -2000, 0);
lc.setLoadFactor(1.0);

% Linear analysis
results = analysis.runLinearAnalysis(lc);
stress  = analysis.performStressAnalysis();

fprintf('Max disp: %.3f mm\n', results.max_displacement * 1e3);
fprintf('Min FOS:  %.2f\n', stress.FOS.min_FOS);

% Visualise
plotter = StructurePlotter(geom, analysis);
plotter.plotDeformedShape();
plotter.plotStressDistribution('combined');
\end{lstlisting}

This example demonstrates that no file import is necessary --- structures can
be defined entirely in code by specifying node coordinates and connectivity
directly.